
\sgaStableTarget{sga6.e10.para.7-11}\textbf{7.11.} For a scheme @@M0@@ over an algebraically closed field @@M1@@, one can define the subgroup of @@M2@@ of elements \emph{algebraically equivalent to zero} as the subgroup of elements of the form
\sgaStableTarget{sga6.e10.eq.7-11-1}@@M3@@
where @@M4@@ and @@M5@@ are two @@M6@@-rational points of a proper, regular, connected algebraic curve @@M7@@ over @@M8@@, @@M9@@ is an element of @@M10@@, and, for every point @@M11@@ of @@M13@@ rational over @@M12@@, one puts
@@M14@@
is the morphism @@M15@@ defined by @@M16@@. The fact that the elements of the form (7.11.1) form a subgroup of @@M17@@ follows by a standard argument, using a product @@M18@@ and its pair of points @@M19@@, @@M20@@, and the fact that through two points of an irreducible quasi-projective algebraic scheme over @@M21@@ (here @@M22@@) one can always pass an irreducible algebraic curve. (A reader unwilling to admit this fact may define algebraic equivalence by the group generated by the elements of the form (7.11.1).) The elements algebraically equivalent to zero form an ideal of @@M23@@, as one sees at once, which is moreover stable under the operations @@M24@@ when these are defined, i.e., when one works with the ``naive'' @@M25@@; this ideal is moreover contained in the augmentation ideal of @@M26@@, as one verifies trivially. One can therefore define a quotient ring
\sgaStableTarget{sga6.e10.eq.7-11-2}@@M27@@
which is even a @@M28@@-ring when one works with naive @@M29@@.

In the same way, one can define a subgroup of elements algebraically equivalent to 0 in @@M30@@, taking @@M31@@ above in @@M32@@, and
